\chapter{Term Rewrite Systems}
\chapterauthor{Jan-Christoph Kassing}

\begin{abstract}
  This chapter describes the syntax and semantics for first-order \emph{Term Rewrite Systems} (TRSs).
  TRSs are the basic formalism behind several categories of the Termination Competition.
  In particular, the categories \emph{TRS Standard}, \emph{TRS Innermost}, \emph{TRS Outermost}, \emph{TRS Relative}, 
  \emph{TRS Derivational Complexity}, \emph{TRS Innermost Derivational Complexity}, \emph{TRS Runtime Complexity}, 
  and \emph{TRS Innermost Runtime Complexity} all use the same underlying notion of first-order rewriting 
  and differ only in the allowed rewrite strategy or in the class of start terms that is considered.
\end{abstract}

\section{Syntax}

\subsection{Grammar}

\begin{tabular}{lll}
trs & ::= & \texttt{(format TRS)} fun* rule* \\
fun & ::= & \texttt{(fun} IDENTIFIER NAT\texttt{)} \\
term & ::= & IDENTIFIER $|$ \texttt{(}IDENTIFIER term+\texttt{)} \\
cost & ::= & \texttt{:cost} NAT \\
rule & ::= & \texttt{(rule} term term cost?\texttt{)} \\
\end{tabular}

\bigskip
Additionally, we have NAT $\in \mathbb{N}$ and IDENTIFIER are intended to be valid SMT-LIB symbols.
In particular, quoted and unquoted identifiers should not be mixed, and reserved words such as 
\texttt{fun}, \texttt{rule}, \texttt{format}, and \texttt{entrypoint} must not be used as simple identifiers.
Function symbols are declared by \texttt{(fun f n)}, where \texttt{n} is the arity of \texttt{f}.
Identifiers that are not declared as function symbols are variables.
The optional cost annotation is part of the general TRS format, but it is only relevant for categories such as relative rewriting.

\textbf{Example.}
The TRS $\mathcal{R} = \{\mathsf{f}(x) \to \mathsf{a}\}$ has the following input syntax:
\begin{center}
\begin{tabular}{l}
\texttt{(format TRS)} \\
\texttt{(fun f 1)} \\
\texttt{(fun a 0)} \\
\texttt{(rule (f x) a)}
\end{tabular}
\end{center}
The TRS $\mathcal{R} = \{\mathsf{f}(x) \to \mathsf{f}(x)\}$ has the following input syntax:
\begin{center}
\begin{tabular}{l}
\texttt{(format TRS)} \\
\texttt{(fun f 1)} \\
\texttt{(rule (f x) (f x))}
\end{tabular}
\end{center}

\subsection{Restrictions}

The grammar above matches the adapted ARI format used on the Termination Portal.
Usually, terms and rules should additionally satisfy the following well-formedness conditions.

\begin{enumerate}
\item Every function symbol is declared exactly once.
\item Function symbols are always used with their declared arity.
\item The left-hand side of a rule is not a variable.
\item Every variable occurring in the right-hand side of a rule also occurs in the corresponding left-hand side.
\end{enumerate}

In the categories \emph{TRS Standard}, \emph{TRS Innermost}, \emph{TRS Outermost}, \emph{TRS Derivational Complexity}, 
\emph{TRS Innermost Derivational Complexity}, \emph{TRS Runtime Complexity}, and \emph{TRS Innermost Runtime Complexity}, rule costs are disallowed.
The category \emph{TRS Relative} uses the same base syntax, but the cost annotations distinguish strict rules from weak rules.
In such relative rules with cost $0$ we omit the last two well-formedness conditions regarding the variables of rules.

\section{Semantics}

We write $\mathcal{T}(\F,\V)$ for the set of all terms over a signature $\F = \biguplus_{k \in \N} \F_k$ of function symbols and a set $\V$ of variables.
More precisely, $\mathcal{T}(\F,\V)$ is the smallest set such that $\V \subseteq \mathcal{T}(\F,\V)$ 
and, whenever $f \in \F_k$ and $t_1,\ldots,t_k \in \mathcal{T}(\F,\V)$, also $f(t_1,\ldots,t_k) \in \mathcal{T}(\F,\V)$.
If $f \in \F_k$, then we say that $f$ has arity $k$.

A \emph{substitution} is a mapping $\sigma : \V \to \mathcal{T}(\F,\V)$ such that $\sigma(x) = x$ for all but finitely many variables $x \in \V$.
Substitutions extend homomorphically to terms, i.e., if $t = f(t_1,\ldots,t_k)$, then $t\sigma = f(t_1\sigma,\ldots,t_k\sigma)$.

A \emph{rewrite rule} is a pair of terms $\ell \to r$ such that $\ell \notin \V$ and $\mathit{Var}(r) \subseteq \mathit{Var}(\ell)$.
A \emph{term rewrite system} is a finite set $\Rules$ of rewrite rules.

\subsection{Positions and Subterms}

Positions are finite sequences of natural numbers, and the empty sequence is denoted by $\varepsilon$.
For a term $t$, the set of its positions is written $\mathit{Pos}(t)$ and is inductively defined as usual: 
we have $\varepsilon \in \mathit{Pos}(t)$ and, if $t = f(t_1,\ldots,t_k)$, then $i.p \in \mathit{Pos}(t)$ 
for every $1 \leq i \leq k$ and every $p \in \mathit{Pos}(t_i)$.
If $p \in \mathit{Pos}(t)$, then $t|_p$ denotes the subterm of $t$ at position $p$, and $t[s]_p$ denotes the term obtained 
from $t$ by replacing the subterm $t|_p$ by $s$.

A term $u$ is a \emph{redex} of $\Rules$ if there exist a rule $\ell \to r \in \Rules$ and a substitution $\sigma$ such that $u = \ell\sigma$.
A term is in \emph{normal form} w.r.t.\ $\Rules$ if it contains no redex.

A \emph{context} is a term with exactly one occurrence of a distinguished hole symbol $\square$.
If $C$ is a context, then $C[t]$ denotes the term obtained by replacing the hole in $C$ by the term $t$.

\subsection{Full Rewriting}

The TRS $\Rules$ induces a one-step rewrite relation ${\to_\Rules} \subseteq \mathcal{T}(\F,\V) \times \mathcal{T}(\F,\V)$.
We define $s \to_\Rules t$ if there exist a position $p \in \mathit{Pos}(s)$, 
a rule $\ell \to r \in \Rules$, and a substitution $\sigma$ 
such that $s|_p = \ell\sigma$ and $t = s[r\sigma]_p$.
Its reflexive-transitive closure is denoted by $\to_\Rules^*$.

The TRS $\Rules$ is \emph{terminating} if $\to_\Rules$ is well founded.

\subsection{Innermost and Outermost Rewriting}

An \emph{innermost rewrite step}, written $s \ito_\Rules t$, is a rewrite step $s \to_\Rules t$ at a position $p$ 
such that all proper subterms of the redex $s|_p$ are in normal form w.r.t.\ $\Rules$.

An \emph{outermost rewrite step}, written $s \oto_{\Rules} t$, is a rewrite step $s \to_\Rules t$ at a position $p$ 
such that there is no redex at a proper prefix of $p$.

The TRS $\Rules$ is \emph{innermost terminating} if if $\ito_\Rules$ is well founded.
It is \emph{outermost terminating} if if $\oto_\Rules$ is well founded.

\subsection{Relative Rewriting}

For two TRSs $\Rules$ and $\mathcal{S}$ over the same signature, 
the \emph{relative rewrite relation} $\to_{\Rules/\mathcal{S}}$ is defined by
\[
  \to_{\Rules/\mathcal{S}} \;=\; \to_{\mathcal{S}}^* \cdot \to_{\Rules} \cdot \to_{\mathcal{S}}^* .
\]
The pair $\Rules/\mathcal{S}$ is \emph{relatively terminating} if $\to_{\Rules/\mathcal{S}}$ is well founded.

In the competition format, relative rewriting can be represented by rule costs, 
where strict rules of $\mathcal{R}$ have cost $1$ and weak rules of $\mathcal{S}$ have cost $0$.
The rules $\ell \to r \in \mathcal{S}$ may ignore the variable conditions 
$\ell \notin \V$ and $\mathit{Var}(r) \subseteq \mathit{Var}(\ell)$.

\section{Complexity Measures}

There are two standard notions of complexity for term rewriting: \emph{derivational complexity} and \emph{runtime complexity}.
For any set $M \subseteq \N$, we write $\sup M$ for its least upper bound, with the convention $\sup \emptyset = 0$.
For a rewrite relation ${\to} \subseteq \mathcal{T}(\F,\V) \times \mathcal{T}(\F,\V)$, 
the \emph{derivation height} of a term $t$ w.r.t.\ $\to$ is
\[
  \mathit{dh}_{\to}(t) = \sup\{ m \in \N \mid \exists t' \in \mathcal{T}(\F,\V) : t \to^{m} t' \}.
\]
Thus, $\mathit{dh}_{\to}(t)$ is the length of the longest $\to$-rewrite sequence starting with $t$.
We measure the \emph{size} $|t|$ of a term $t$ 
by the number of occurrences of function symbols and variables in $t$,
i.e., $|x| = 1$ for variables $x \in \V$ and
\[
  |f(t_1,\ldots,t_k)| = 1 + \sum_{i=1}^{k} |t_i|.
\]

\subsection{Derivational Complexity}

The \emph{derivational complexity} of a rewrite relation $\to$ is the function
\[
  \mathit{dc}_{\to}(n) = \sup\{ \mathit{dh}_{\to}(t) \mid t \in \mathcal{T}(\F,\V),\ |t| \leq n \}.
\]
Hence, $\mathit{dc}_{\to}(n)$ is the maximal length of a $\to$-rewrite sequence 
starting with an arbitrary term of size at most $n$.

For the TRS $\Rules$, its full derivational complexity is $\mathit{dc}_{\to_\Rules}$.
Analogously, the \emph{innermost derivational complexity} of $\Rules$ is obtained 
by replacing $\to_\Rules$ with $\ito_\Rules$.

\subsection{Runtime Complexity}

In contrast, runtime complexity restricts the start terms to \emph{basic terms}.
To this end, we decompose the signature as $\F = \mathcal{C} \uplus \mathcal{D}$, 
where $f \in \mathcal{D}$ if $f$ occurs at the root of the left-hand side of some rule in $\Rules$.
The symbols in $\mathcal{C}$ and $\mathcal{D}$ are called \emph{constructors} and \emph{defined symbols}, respectively.

A term $t \in \mathcal{T}(\F,\V)$ is \emph{basic} if $t = f(t_1,\ldots,t_k)$ 
for some defined symbol $f \in \mathcal{D}$ and constructor terms $t_1,\ldots,t_k \in \mathcal{T}(\mathcal{C},\V)$.
We write $\mathcal{T}_{\mathit{B}}^{\mathcal{R}}$ for the set of all basic terms w.r.t.\ $\mathcal{R}$.

The \emph{runtime complexity} of a rewrite relation $\to$ w.r.t.\ $\Rules$ is the function
\[
  \mathit{rc}_{\to,\Rules}(n) = \sup\{ \mathit{dh}_{\to}(t) \mid t \in \mathcal{T}_{\mathit{B}}^{\mathcal{R}},\ |t| \leq n \}.
\]
Thus, runtime complexity measures the maximal length of a rewrite sequence starting with a basic term of size at most $n$.
For the TRS $\Rules$, its full runtime complexity is $\mathit{rc}_{\to_\Rules,\Rules}$, 
and its \emph{innermost runtime complexity} is defined analogously using $\ito_\Rules$ instead of $\to_\Rules$.

\section{Competition Categories}

The TRS formalism from this chapter is used directly by the following categories.

\subsection{TRS Standard}

\textbf{Motivation.}
This is the basic termination problem for unrestricted first-order rewriting.
Many other rewriting categories are refinements or extensions of this setting, so it serves as a common baseline for termination techniques.

\begin{formalproblem}
  {Formal Problem Description}
  {A TRS $\Rules$}
  {Does $\Rules$ terminate?}
  {$\texttt{YES} \quad | \quad \texttt{NO} \quad | \quad \texttt{MAYBE}$}
\end{formalproblem}

\textbf{Example.}
The TRS $\mathcal{R} = \{\mathsf{f}(x) \to \mathsf{a}\}$ has the following input syntax:
\begin{center}
\begin{tabular}{l}
\texttt{(format TRS)} \\
\texttt{(fun f 1)} \\
\texttt{(fun a 0)} \\
\texttt{(rule (f x) a)}
\end{tabular}
\end{center}
Note that $\mathcal{R}$ is terminating, and thus, the expected output should be \texttt{YES}.
The TRS $\mathcal{R} = \{\mathsf{f}(x) \to \mathsf{f}(x)\}$ has the following input syntax:
\begin{center}
\begin{tabular}{l}
\texttt{(format TRS)} \\
\texttt{(fun f 1)} \\
\texttt{(rule (f x) (f x))}
\end{tabular}
\end{center}
Note that $\mathcal{R}$ is not terminating, and thus, the expected output should be \texttt{NO}.

\subsection{TRS Relative}

\textbf{Motivation.}
Relative termination isolates the rules whose repeated use should be ruled out while allowing auxiliary rules to be used freely in between.
This is useful for modular proofs and transformations that reduce a problem to showing that a strict part terminates relative to a weak part.

\begin{formalproblem}
  {Formal Problem Description}
  {A relative TRS $\Rules/\mathcal{S}$}
  {Is $\Rules/\mathcal{S}$ relatively terminating?}
  {$\texttt{YES} \quad | \quad \texttt{NO} \quad | \quad \texttt{MAYBE}$}
\end{formalproblem}
\textbf{Example.}
The relative TRS $\mathcal{R}/\mathcal{S}$ with strict part $\mathcal{R} = \{\mathsf{a} \to \mathsf{b}\}$ 
and weak part $\mathcal{S} = \{\mathsf{b} \to \mathsf{b}\}$ has the following input syntax:
\begin{center}
\begin{tabular}{l}
\texttt{(format TRS)} \\
\texttt{(fun a 0)} \\
\texttt{(fun b 0)} \\
\texttt{(rule a b :cost 1)} \\
\texttt{(rule b b :cost 0)}
\end{tabular}
\end{center}
There is at most one strict step, and thus, the expected output should be \texttt{YES}.
The relative TRS $\mathcal{R}/\mathcal{S}$ with strict part $\mathcal{R} = \{\mathsf{a} \to \mathsf{b}\}$ 
and weak part $\mathcal{S} = \{\mathsf{b} \to \mathsf{a}\}$ has the following input syntax:
\begin{center}
\begin{tabular}{l}
\texttt{(format TRS)} \\
\texttt{(fun a 0)} \\
\texttt{(fun b 0)} \\
\texttt{(rule a b :cost 1)} \\
\texttt{(rule b a :cost 0)}
\end{tabular}
\end{center}
It yields an infinite sequence using the strict rule infinitely often, and thus, the expected output should be \texttt{NO}.

\subsection{TRS Innermost}

\textbf{Motivation.}
This category studies termination under innermost rewriting, which closely corresponds to call-by-value evaluation.
Since many programming languages evaluate arguments before calling a function, 
innermost termination is an especially important restricted strategy.

\begin{formalproblem}
  {Formal Problem Description}
  {A TRS $\Rules$}
  {Does $\Rules$ innermost terminate?}
  {$\texttt{YES} \quad | \quad \texttt{NO} \quad | \quad \texttt{MAYBE}$}
\end{formalproblem}

\textbf{Example.}
The TRS $\mathcal{R} = \{\mathsf{f}(x) \to x\}$ has the following input syntax:
\begin{center}
\begin{tabular}{l}
\texttt{(format TRS)} \\
\texttt{(fun f 1)} \\
\texttt{(fun a 0)} \\
\texttt{(rule (f x) x)}
\end{tabular}
\end{center}
Every innermost step strictly decreases the term size, and thus, the expected output should be \texttt{YES}.
The TRS $\mathcal{R} = \{\mathsf{g}(x) \to \mathsf{g}(x)\}$ has the following input syntax:
\begin{center}
\begin{tabular}{l}
\texttt{(format TRS)} \\
\texttt{(fun g 1)} \\
\texttt{(rule (g x) (g x))}
\end{tabular}
\end{center}
Note that $\mathcal{R}$ is not innermost terminating since $\mathsf{g}(x)$ rewrites to itself by an innermost step, 
and thus, the expected output should be \texttt{NO}.

\subsection{TRS Outermost}

\textbf{Motivation.}
This category studies termination under outermost rewriting, 
which is closely related to call-by-name or lazy evaluation.
Although more restrictive than full rewriting, this strategy is important for languages 
and semantics that evaluate the outermost call first.

\begin{formalproblem}
  {Formal Problem Description}
  {A TRS $\Rules$}
  {Does $\Rules$ outermost terminate?}
  {$\texttt{YES} \quad | \quad \texttt{NO} \quad | \quad \texttt{MAYBE}$}
\end{formalproblem}

\textbf{Example.}
The TRS $\mathcal{R} = \{\mathsf{f}(x) \to \mathsf{a}\}$ has the following input syntax:
\begin{center}
\begin{tabular}{l}
\texttt{(format TRS)} \\
\texttt{(fun f 1)} \\
\texttt{(fun a 0)} \\
\texttt{(rule (f x) a)}
\end{tabular}
\end{center}
Every outermost redex reduces to the normal form $\mathsf{a}$, and thus, the expected output should be \texttt{YES}.
The TRS $\mathcal{R} = \{\mathsf{f}(x) \to \mathsf{f}(x)\}$ has the following input syntax:
\begin{center}
\begin{tabular}{l}
\texttt{(format TRS)} \\
\texttt{(fun f 1)} \\
\texttt{(rule (f x) (f x))}
\end{tabular}
\end{center}
Note that $\mathcal{R}$ is not outermost terminating since the outermost redex rewrites to itself forever, and thus, the expected output should be \texttt{NO}.

\subsection{TRS Derivational Complexity}

\textbf{Motivation.}
This category asks not only whether rewriting terminates, 
but how long rewrite sequences can become as a function of the input size.
It measures the worst-case asymptotic complexity of unrestricted rewriting on arbitrary start terms.

\begin{formalproblem}
  {Formal Problem Description}
  {A TRS $\Rules$}
  {What is the asymptotic derivational complexity of $\Rules$?}
  {\texttt{WORST\_CASE(}$\ell\texttt{,}$u\texttt{)}, where}
  \[
    \ell,u \in \{\mathrm{Pol}_0, \mathrm{Pol}_1, \mathrm{Pol}_2, \ldots, \mathrm{Exp}, \mathrm{DExp}, \omega\}
  \]
\end{formalproblem}

\textbf{Example.}
The TRS $\mathcal{R} = \{\mathsf{f}(x) \to \mathsf{a}\}$ has the following input syntax:
\begin{center}
\begin{tabular}{l}
\texttt{(format TRS)} \\
\texttt{(fun f 1)} \\
\texttt{(fun a 0)} \\
\texttt{(rule (f x) a)}
\end{tabular}
\end{center}
Here, a rewrite sequence starting with $\mathsf{f}^n(x)$ has length $n$ in the worst-case, 
so the derivational complexity has class $\mathrm{Pol}_1$, and the expected output can be \texttt{WORST\_CASE(}$\mathrm{Pol}_1$\texttt{,}$\mathrm{Pol}_1$\texttt{)}.
The TRS $\mathcal{R} = \{\mathsf{f}(x) \to \mathsf{f}(x)\}$ has the following input syntax:
\begin{center}
\begin{tabular}{l}
\texttt{(format TRS)} \\
\texttt{(fun f 1)} \\
\texttt{(rule (f x) (f x))}
\end{tabular}
\end{center}
For this system, derivation heights are infinite, so the derivational complexity has class $\omega$, and the expected output can be \texttt{WORST\_CASE(}$\omega$\texttt{,}$\omega$\texttt{)}.

\subsection{TRS Innermost Derivational Complexity}

\textbf{Motivation.}
This category measures the asymptotic complexity of a TRS 
when only innermost rewrite sequences are considered.
It is especially relevant when the intended execution model follows call-by-value evaluation.

\begin{formalproblem}
  {Formal Problem Description}
  {A TRS $\Rules$}
  {What is the asymptotic innermost derivational complexity of $\Rules$?}
  {\texttt{WORST\_CASE(}$\ell\texttt{,}$u\texttt{)}, where}
  \[
    \ell,u \in \{\mathrm{Pol}_0, \mathrm{Pol}_1, \mathrm{Pol}_2, \ldots, \mathrm{Exp}, \mathrm{DExp}, \omega\}
  \]
\end{formalproblem}

\textbf{Example.}
The TRS $\mathcal{R} = \{\mathsf{f}(x) \to \mathsf{a}\}$ has the following input syntax:
\begin{center}
\begin{tabular}{l}
\texttt{(format TRS)} \\
\texttt{(fun f 1)} \\
\texttt{(fun a 0)} \\
\texttt{(rule (f x) a)}
\end{tabular}
\end{center}
Again, an innermost rewrite sequence starting with $\mathsf{f}^n(x)$ has length $n$ in the worst-case, so the innermost derivational complexity has class $\mathrm{Pol}_1$, and the expected output can be \texttt{WORST\_CASE(}$\mathrm{Pol}_1$\texttt{,}$\mathrm{Pol}_1$\texttt{)}.
The TRS $\mathcal{R} = \{\mathsf{g}(x) \to \mathsf{g}(x)\}$ has the following input syntax:
\begin{center}
\begin{tabular}{l}
\texttt{(format TRS)} \\
\texttt{(fun g 1)} \\
\texttt{(rule (g x) (g x))}
\end{tabular}
\end{center}
For this system, innermost derivation heights are infinite, so the innermost derivational complexity has class $\omega$, and the expected output can be \texttt{WORST\_CASE(}$\omega$\texttt{,}$\omega$\texttt{)}.

\subsection{TRS Runtime Complexity}

\textbf{Motivation.}
Runtime complexity focuses on terms that represent an algorithm symbol applied to already constructed input data.
This better matches the usual notion of program runtime than derivational complexity on arbitrary terms.

\begin{formalproblem}
  {Formal Problem Description}
  {A TRS $\Rules$}
  {What is the asymptotic runtime complexity of $\Rules$?}
  {\texttt{WORST\_CASE(}$\ell\texttt{,}$u\texttt{)}, where}
  \[
    \ell,u \in \{\mathrm{Pol}_0, \mathrm{Pol}_1, \mathrm{Pol}_2, \ldots, \mathrm{Exp}, \mathrm{DExp}, \omega\}
  \]
\end{formalproblem}

\textbf{Example.}
The TRS $\mathcal{R} = \{\mathsf{f}(\mathsf{z}) \to \mathsf{z},\; \mathsf{f}(\mathsf{s}(x)) \to \mathsf{f}(x)\}$ has the following input syntax:
\begin{center}
\begin{tabular}{l}
\texttt{(format TRS)} \\
\texttt{(fun f 1)} \\
\texttt{(fun s 1)} \\
\texttt{(fun z 0)} \\
\texttt{(rule (f z) z)} \\
\texttt{(rule (f (s x)) (f x))}
\end{tabular}
\end{center}
Starting from a basic term of the form $\mathsf{f}(\mathsf{s}^n(\mathsf{z}))$, 
the number of steps is linear in the input size, so the runtime complexity has class $\mathrm{Pol}_1$, 
and the expected output can be \texttt{WORST\_CASE(}$\mathrm{Pol}_1$\texttt{,}$\mathrm{Pol}_1$\texttt{)}.
The TRS $\mathcal{R} = \{\mathsf{f}(x) \to \mathsf{f}(x)\}$ has the following input syntax:
\begin{center}
\begin{tabular}{l}
\texttt{(format TRS)} \\
\texttt{(fun f 1)} \\
\texttt{(rule (f x) (f x))}
\end{tabular}
\end{center}
Since $\mathsf{f}(c)$ is a basic term for every constructor term $c$,
the runtime complexity has class $\omega$, 
and the expected output can be \texttt{WORST\_CASE(}$\omega$\texttt{,}$\omega$\texttt{)}.

\subsection{TRS Innermost Runtime Complexity}

\textbf{Motivation.}
This category combines the restriction to basic start terms with the restriction to innermost rewriting.
It captures the asymptotic cost of call-by-value computations on TRSs that represent algorithms on fixed input data.

\begin{formalproblem}
  {Formal Problem Description}
  {A TRS $\Rules$}
  {What is the asymptotic innermost runtime complexity of $\Rules$?}
  {\texttt{WORST\_CASE(}$\ell\texttt{,}$u\texttt{)}, where}
  \[
    \ell,u \in \{\mathrm{Pol}_0, \mathrm{Pol}_1, \mathrm{Pol}_2, \ldots, \mathrm{Exp}, \mathrm{DExp}, \omega\}
  \]
\end{formalproblem}

\textbf{Example.}
The TRS $\mathcal{R} = \{\mathsf{f}(\mathsf{z}) \to \mathsf{z},\; \mathsf{f}(\mathsf{s}(x)) \to \mathsf{f}(x)\}$ has the following input syntax:
\begin{center}
\begin{tabular}{l}
\texttt{(format TRS)} \\
\texttt{(fun f 1)} \\
\texttt{(fun s 1)} \\
\texttt{(fun z 0)} \\
\texttt{(rule (f z) z)} \\
\texttt{(rule (f (s x)) (f x))}
\end{tabular}
\end{center}
On basic terms of the form $\mathsf{f}(\mathsf{s}^n(\mathsf{z}))$, 
each innermost step removes one constructor $\mathsf{s}$, so the innermost runtime complexity has class $\mathrm{Pol}_1$, 
and the expected output can be \texttt{WORST\_CASE(}$\mathrm{Pol}_1$\texttt{,}$\mathrm{Pol}_1$\texttt{)}.
The TRS $\mathcal{R} = \{\mathsf{f}(x) \to \mathsf{f}(x)\}$ has the following input syntax:
\begin{center}
\begin{tabular}{l}
\texttt{(format TRS)} \\
\texttt{(fun f 1)} \\
\texttt{(rule (f x) (f x))}
\end{tabular}
\end{center}
Note that every basic term rooted by $\mathsf{f}$ rewrites to itself by an innermost step, so the innermost runtime complexity has class $\omega$ and the expected output can be \texttt{WORST\_CASE(}$\omega$\texttt{,}$\omega$\texttt{)}.